\documentclass[10pt]{amsart}
\usepackage{geometry}                % See geometry.pdf to learn the layout options. There are lots.
\geometry{letterpaper}                   % ... or a4paper or a5paper or ... 
%\geometry{landscape}                % Activate for for rotated page geometry
%\usepackage[parfill]{parskip}    % Activate to begin paragraphs with an empty line rather than an indent
\usepackage{graphicx}
\usepackage{amssymb}
\usepackage{epstopdf}
\usepackage{natbib}
\DeclareGraphicsRule{.tif}{png}{.png}{`convert #1 `dirname #1`/`basename #1 .tif`.png}

\title{Some Solutions to the Planetary Geostrophic Equations}
\author{TWNH}
\date{\today}                                           % Activate to display a given date or no date

\begin{document}
\maketitle

\citet{peterson&callies26} (see also \citealt{peterson&callies23}) consider how bottom-intensified vertical mixing induces a basin-scale circulation in various simple non-flat-bottomed basins.
They write equations for a bottom Ekman layer and couple them to an equation for the geostrophic contour vorticity balance, in the parlance of \citet{waldman&giordani23}, which is a type of barotropic vorticity equation.

\section{Governing Equations}
The flow satisfies the (non-dimensional) planetary-geostrophic equations:
\begin{align}
- f v & = - \frac{\partial p}{\partial x} + \epsilon^2 \frac{\partial}{\partial z} \left( \nu \frac{\partial u}{\partial z} \right), \label{eqn:u_mom} \\
f  u & = - \frac{\partial p}{\partial y} + \epsilon^2 \frac{\partial}{\partial z} \left( \nu \frac{\partial v}{\partial z} \right),  \label{eqn:v_mom} \\
\end{align}
(see Eqs.~(9)--(11) of \citealt{peterson&callies26}, Eq.~(3.26) of \citealt{klinger&haine19}, and Eq.~(2.274) of \citealt{vallis06}).
The boundary conditions are no flow at the bottom and a specified stress at the surface:
\begin{align}
\text{Surface~}z = 0:
\epsilon^2 \nu \partial_z (u, v) & = (\tau^x , \tau^y ), 
 \\
\text{Bottom~}z = -H:
(u, v) & = 0 ,
\end{align}
The flow is non-divergent:
\begin{align}
\nabla \cdot {\mathbf u} & = 0 ,
\end{align}
the buoyancy field is hydrostatic:
\begin{align}
\frac{\partial p}{\partial z} & = b ,
\end{align}
and the buoyancy field satisfies the advection-diffusion-relaxation equation
\begin{align}
\mu \rho \left( \frac{\partial b}{\partial t} + {\mathbf u} \cdot \nabla b \right) = \epsilon^2 \frac{\partial }{\partial z} \left( \kappa \frac{\partial b}{\partial z} \right) + \gamma \left( B - b \right).
\label{eqn:b}
\end{align}
(see Eqs.~(11)--(13) of \citealt{peterson&callies26}).
The boundary conditions on buoyancy are:
\begin{align}
\text{Surface~}z = 0:
b &= 0
\label{eqn:sfc_b_bc}
 \\
\text{Bottom~}z = -H:
\frac{\partial b}{\partial z} &= 0 ,
\label{eqn:bot_b_bc}
\end{align}
and the initial condition is $b = z$.
The buoyancy field relaxes to a specified field $B(x, y, z, t)$ with rate $\gamma$.

The diffusivity and viscosity are given by
\begin{align}
\kappa(z) & = \kappa_0 + \exp \left( - \alpha (z + H) \right) , \label{eqn:kappa_def} \\
\nu & = \text{constant}
\end{align}
\citet{peterson&callies26} Eq.~(17).

\section{Vertically-Integrated Flow}
The vertically-integrated flow $(U, V)$ is given by:
\begin{align}
U & \equiv \int_{-H}^{0} u \; dz = - \int_{-H}^0 z \partial_z u \; dz , \nonumber \\
V & \equiv \int_{-H}^{0} v \; dz = - \int_{-H}^0 z \partial_z v \; dz ,
\label{eqn:UV_defs}
\end{align}
exploiting the bottom boundary condition on $(u, v)$.
The equations for $(U, V)$ come from vertically integrating the momentum equations (\ref{eqn:u_mom}), (\ref{eqn:v_mom}):
\begin{align}
- f V & = -  \int_{-H}^{0} \frac{\partial p}{\partial x} \; dz + \epsilon^2   \int_{-H}^{0} \frac{\partial}{\partial z} \left( \nu \frac{\partial u}{\partial z} \right) \; dz \label{eqn:U_mom0} \\
f  U & = -  \int_{-H}^{0} \frac{\partial p}{\partial y} \; dz + \epsilon^2  \int_{-H}^{0} \frac{\partial}{\partial z} \left( \nu \frac{\partial v}{\partial z} \right) \; dz . \label{eqn:V_mom0}
\end{align}
The depth-integrated pressure gradient in the $x$ direction is:
\begin{align}
 \int_{-H}^{0} \frac{\partial p}{\partial x} \; dz & = - \frac{\partial H}{\partial x} p(x,z = -H) +  \frac{\partial }{\partial x} \int_{-H}^0 p \; dz , \\
 & = - p_b \frac{\partial H}{\partial x} + \frac{\partial }{\partial x} \left[ p_b H - \int_{-H}^0 z \frac{\partial p}{\partial z}  \right], \\
  & = H \frac{\partial p_b}{\partial x} - \frac{\partial }{\partial x} \int_{-H}^0 z b \; dz ,
 \end{align}
 using the Liebniz integration formula\footnote{The Liebniz integration formula reads:
 \begin{align}
 \partial_x \left( \int_{a(x)}^{b(x)} f(x,y) \; dy \right) = \partial_x b f(x,b(x)) - \partial_x a f(x,a(x)) +  \int_{a(x)}^{b(x)} \partial_x f(x, y) \; dy . \nonumber
 \end{align}
 }, where $p_b = p(x,-H(x))$ is the bottom pressure, and integration by parts\footnote{
 Integration by parts gives:
 \begin{align}
\int_{-H}^0 z \frac{\partial p}{\partial z} \; dz & = \left[ z p \right]_{-H}^0 -   \int_{-H}^0 p \; dz , \nonumber \\ 
& = p_b H -  \int_{-H}^0 p \; dz . \nonumber
 \end{align}
 }.
 The viscous term in the $x$ direction is:
 \begin{align}
  \int_{-H}^{0} \frac{\partial}{\partial z} \left( \nu \frac{\partial u}{\partial z} \right) \; dz & = \left[  \nu \frac{\partial u}{\partial z}  \right]_{-H}^{0}  = \frac{\tau^x}{ \epsilon^2}  - \left. \nu \frac{\partial u}{\partial z} \right|_{z = -H} ,
 \end{align}
which uses the boundary condition (\ref{eqn:sfc_bcs}).
  
 Therefore (\ref{eqn:U_mom0}) and (\ref{eqn:V_mom0}) lead to:
 \begin{align}
- f V & =  - H \frac{\partial p_b}{\partial x} - \frac{\partial \chi}{\partial x}  - \tau^x_b + \tau^x , \label{eqn:U_mom} \\
f  U & = - H \frac{\partial p_b}{\partial y} - \frac{\partial \chi }{\partial y}  - \tau^y_b + \tau^y \label{eqn:V_mom} ,
\end{align}
where
\begin{align}
\chi & = - \int_{-H}^0 z b \; dz
\label{eqn:baroclinic_PE}
\end{align}
is the baroclinic potential energy, and $\epsilon^2  \nu \left. \partial_z (u,v) \right|_{z = -H}  \equiv (\tau^x_b, \tau^y_b) = \tau_b$ is the bottom stress on the fluid.

Next, divide by water depth $H$:
 \begin{align}
- \frac{f V}{H} & =  - \frac{\partial p_b}{\partial x} - \frac{1}{H} \frac{\partial \chi}{\partial x} + \frac{\tau^x - \tau^x_b}{H}   \\
\frac{f  U}{H} & = - \frac{\partial p_b}{\partial y} - \frac{1}{H} \frac{\partial \chi }{\partial y} + \frac{\tau^y - \tau^y_b}{H} ,
\end{align}
and take the curl\footnote{
Recall that 
\begin{align}
{\hat{\mathbf z}} \cdot \left( \nabla \times {\mathbf U} \right) & = \partial_x V - \partial_y U .
\end{align}
}
 \begin{align}
\frac{\partial}{\partial y} \left\{- \frac{f V}{H} \right\}  - \frac{\partial}{\partial x} \left\{\frac{f  U}{H} \right\} & = 
\frac{\partial}{\partial y} \left\{ - \frac{\partial p_b}{\partial x} - \frac{1}{H} \frac{\partial \chi}{\partial x}  + \frac{\tau^x - \tau^x_b}{H} \right\} \nonumber. \\
& -  \frac{\partial}{\partial x} \left\{ - \frac{\partial p_b}{\partial y} - \frac{1}{H} \frac{\partial \chi }{\partial y}  + \frac{\tau^y - \tau^y_b}{H} \right\}, \\
\implies
\frac{\partial}{\partial y} \left\{- \frac{f }{H} \frac{\partial \Psi}{\partial x} \right\}  + \frac{\partial}{\partial x} \left\{\frac{f }{H} \frac{\partial \Psi}{\partial y} \right\} & = 
- \frac{\partial}{\partial y} \left\{ \frac{1}{H} \frac{\partial \chi}{\partial x} \right\} + \frac{\partial}{\partial x} \left\{ \frac{1}{H} \frac{\partial \chi }{\partial y}  \right\}  - {\hat{\mathbf z}} \cdot \left( \nabla \times \frac{{\mathbf \tau - \tau_b}}{H} \right), \nonumber \\
\implies
- \frac{\partial}{\partial y} \left(\frac{f }{H} \right) \, \frac{\partial \Psi}{\partial x}  + \frac{\partial}{\partial x} \, \left( \frac{f }{H}  \right) \frac{\partial \Psi}{\partial y} & = 
- \frac{\partial}{\partial y} \left( \frac{1}{H} \right) \, \frac{\partial \chi}{\partial x} + \frac{\partial}{\partial x} \, \left( \frac{1}{H} \right) \frac{\partial \chi }{\partial y}   - {\hat{\mathbf z}} \cdot \left( \nabla \times \frac{{\mathbf \tau - \tau_b}}{H} \right), \label{eqn:GCVB0} \\
\implies
J\left( \frac{f}{H} , \Psi \right)  & =  J \left( \frac{1}{H}, \chi \right)
- {\hat{\mathbf z}} \cdot \left( \nabla \times \frac{{\mathbf \tau - \tau_b}}{H} \right) ,
\label{eqn:GCVB}
\end{align}
for streamfunction of the depth-integrated horizontal flow $\Psi$ where
\begin{align}
V & = \phantom{-} \frac{\partial \Psi}{\partial x}, \nonumber \\
U & = - \frac{\partial \Psi}{\partial y} ,
\end{align}
and Jacobian $J(A,B) = \partial_x A \, \partial_y B - \partial_y A \, \partial_x B$.
Notice how this recipe (vertically integrate, divide by $H$, then take the curl) eliminates the bottom pressure.
The final result (\ref{eqn:GCVB}) is the same as Eq.~(20) of \citet{peterson&callies26}.
It is called (a simplified form of) the \textbf{geostrophic contour vorticity balance} (GCVB) by \citet{waldman&giordani23}, although \citet{peterson&callies26} call it the ``barotropic vorticity equation''.
It is not the same equation as used by \citet{nost&isachsen03} (which \citet{waldman&giordani23} call the ``transport divergence vorticity balance'').

Notice how the GCVB (\ref{eqn:GCVB}) is an equation for the streamfunction of the depth-integrated flow $\Psi$ given the buoyancy field (and hence the JEBAR term) and the surface and bottom stresses.
The GCVB shows how depth-integrated flow $\Psi$ across $f/H$ contours is driven by (i) the joint effect of baroclinicity and relief (JEBAR), (ii) curl of wind stress divided by $H$, (iii) curl of the bottom stress divided by $H$.

\section{Solution to Frictional Geostrophic Equations in Steady State}
Start from (\ref{eqn:u_mom})--(\ref{eqn:v_mom}) and write the pressure as
\begin{align}
p(x,y,z) & = p_s(x,y) + \int_{z}^{0} b(x,y,z') \; dz' , \label{eqn:hydrostatic_p} \\
\implies \partial_x p(x,y,z) & = \partial_x p_s (x, y) + \int_{z}^{0} \partial_x b(x,y,z') \; dz'  \label{eqn:hydrostatic_p_pgf} \\
\implies p_b(x, y) & \equiv p(x, y, z=-H(x,y)) = p_s(x,y) + \int_{-H(x,y)}^{0} b(x,y,z') \; dz' , \label{eqn:pb} 
\end{align}
where $p_s(x,y)$ is the surface pressure (and assuming that $p_s \ll p(-H)$).
Solve these frictional geostrophic equations using a Green's function, defined so that $\mathfrak{u} = u + i v$, thus
\begin{align}
\epsilon^2 \frac{\partial}{\partial z} \left( \nu  \frac{\partial \mathfrak{u}}{\partial z} \right)  - i f  \mathfrak{u} & = \mathfrak{f} ,  \\
\text{Surface~}z = 0:
\epsilon^2 \nu \partial_z \mathfrak{u} & = \tau^x + i \tau^y , \label{eqn:sfc_bcs}
 \\
\text{Bottom~}z = -H:
\mathfrak{u} & = 0 ,
\end{align}
where $\mathfrak{f} = \frac{\partial p}{\partial x} +  i \frac{\partial p}{\partial y}$ is the pressure gradient from (\ref{eqn:hydrostatic_p_pgf}).

The corresponding Green's function satisfies:
\begin{align}
\epsilon^2 \frac{\partial}{\partial z} \left( \nu  \frac{\partial G_\mathfrak{u}}{\partial z} \right)  - i f  G_\mathfrak{u} & = \delta(z - \xi) ,  \\
\text{Surface~}z = 0:
\epsilon^2 \nu \partial_z G_\mathfrak{u} & = 0 ,
 \\
\text{Bottom~}z = -H:
G_\mathfrak{u} & = 0 .
\end{align}
We find that the solution for constant viscosity $\nu = \nu_0$ is:  
\begin{align}
 G_{ \mathfrak{u}}(z;\xi) & = 
\frac{ i^{3/2} \mathrm{e}^{- \sqrt{i} \varphi (z + \xi)} }{2  \nu_0 \epsilon^2 \varphi  \left(1 + \mathrm{e}^{2 \sqrt{i} \varphi H \left(x \right)} \right)}
\begin{cases} 
\displaystyle - \mathrm{e}^{2 \sqrt{i} \varphi \xi} + \mathrm{e}^{2 \sqrt{i} \varphi \left( z + H(x)\right)} + \mathrm{e}^{2 \sqrt{i} \varphi \left( z + \xi + H(x)\right)} - 1
 & \text{for}\: z \leq \xi , \\
\displaystyle - \mathrm{e}^{2 \sqrt{i} \varphi z} + \mathrm{e}^{2 \sqrt{i} \varphi \left( \xi + H(x)\right)} + \mathrm{e}^{2 \sqrt{i} \varphi \left( z + \xi + H(x)\right)} - 1
 & \text{for}\: z \geq \xi ,
\end{cases} , \label{eqn:FG_Gsfn} \\
\implies  G_{ \mathfrak{u}}(z;0) & = 
\frac{ i^{3/2} \mathrm{e}^{- \sqrt{i} \varphi z } }{ \nu_0 \epsilon^2 \varphi  \left(1 + \mathrm{e}^{2 \sqrt{i} \varphi H \left(x \right)} \right))}
\left[ \mathrm{e}^{2 \sqrt{i} \varphi ( z  + H{\left(x \right))}} - 1 \right] , \nonumber \\
\implies \partial_z \left. G_{ \mathfrak{u}}(z;\xi) \right|_{z = -H(x)} & = 
\frac{- \mathrm{e}^{-\sqrt{i} \varphi (\xi - H(x))  } }{ \nu_0 \epsilon^2 \left(1 + \mathrm{e}^{2 \sqrt{i} \varphi H \left(x \right)} \right)}
\left( 1 + \mathrm{e}^{2 \sqrt{i} \varphi \xi}\right) , \label{eqn:FG_Gsfn_bottom_deriv} \\
\implies \partial_z \left. G_{ \mathfrak{u}}(z;0) \right|_{z = -H(x)} & = 
\frac{- 2 \mathrm{e}^{\sqrt{i} \varphi H(x)  } }{ \nu_0 \epsilon^2 \left(1 + \mathrm{e}^{2 \sqrt{-i} \varphi H \left(x \right)} \right)}
 , \label{eqn:FG_Gsfn_bottom_deriv_top}
\end{align}
where $\varphi = \sqrt{\left( f / \nu_0 \right) } / \epsilon$ (see \texttt{Example\_barotropic\_GeostrophicFlow\_v0.2.ipynb}).
And then 
\begin{align}
\mathfrak{u} (z;x) & = \int_{-H(x)}^{0} G_{ \mathfrak{u}} (z; \xi) \left[ \partial_x p( x, \xi) + i \partial_y p( x, \xi) \right] \; d \xi 
 -   G_{\mathfrak{u}} (z; 0)  \left( \tau^x + i \tau^y \right)  ,
\end{align}
The baroclinic pressure (due to $b$) adds to the barotropic pressure (due to $p_s(x,y)$) in the integrand of this expression. 
Otherwise, everything else is known (remember that $x$ in these expressions stands for $(x,y)$ in general).

Also, the bottom current gradient is given by
\begin{align}
-\partial_z \left. \mathfrak{u} (z;x) \right|_{z = -H(x)} & =  \int_{-H(x)}^{0} 
 \frac{ \mathrm{e}^{-\sqrt{i} \varphi (\xi - H(x))  } }{ \nu_0 \epsilon^2 \left(1 + \mathrm{e}^{2 \sqrt{i} \varphi H \left(x \right)} \right)}
\left( 1 + \mathrm{e}^{2 \sqrt{i} \varphi \xi}\right) 
 \left[ \partial_x p( x, \xi) + i \partial_y p( x, \xi) \right] \; d \xi \nonumber \\
& - \frac{ 2 \mathrm{e}^{\sqrt{i} \varphi H(x)  } }{ \nu_0 \epsilon^2 \left(1 + \mathrm{e}^{2 \sqrt{i} \varphi H \left(x \right)} \right)}
\left( \tau^x + i \tau^y \right) ,
\end{align}
using (\ref{eqn:FG_Gsfn_bottom_deriv}) and (\ref{eqn:FG_Gsfn_bottom_deriv_top}), and therefore the bottom stress $\tau_b$ on the fluid is
\begin{align}
\tau_b = - \nu_0 \epsilon^2 \partial_z \left. \mathfrak{u} (z;x) \right|_{z = -H(x)} & = \int_{-H(x)}^{0} 
\frac{ \mathrm{e}^{-\sqrt{i} \varphi (\xi - H(x))  } }{ 1 + \mathrm{e}^{2 \sqrt{i} \varphi H \left(x \right)} }
\left( 1 + \mathrm{e}^{2 \sqrt{i} \varphi \xi}\right) 
 \left[ \partial_x p( x, \xi) + i \partial_y p( x, \xi) \right] \; d \xi \nonumber \\
&-\frac{ 2 \mathrm{e}^{\sqrt{i} \varphi H(x)  } }{ 1 + \mathrm{e}^{2 \sqrt{i} \varphi H \left(x \right)} }
 \left( \tau^x + i \tau^y \right) .
 \label{eqn:bottom_stress_expression}
 \end{align}
Notice how the derivatives of the surface pressure part of $p$ end up multiplying expressions that are functions of $x$.

\section{Frictional Geostrophic Equations Closure with the GCVB Equation}

The problem is closed using the GCVB equation, as follows (\citealt{salmon98}, p.~170).
This determines the unknown surface pressure field given the surface stress (and vice versa).
In particular, work from (\ref{eqn:U_mom})--(\ref{eqn:V_mom}),
\begin{align}
\mathfrak{U} (z;x) & \equiv U + i V = \int_{-H(x)}^{0} \mathfrak{u} (z;x) \; dz , \label{eqn:depth_integrated_flow_def} \\
& = \int_{-H(x)}^0 \left\{ \int_{-H(x)}^{0} G_{ \mathfrak{u}} (z; \xi) \left[ \partial_x p( x, \xi) + i \partial_y p( x, \xi) \right] \; d \xi 
-  G_{\mathfrak{u}} (z; 0)  \left( \tau^x + i \tau^y \right)  \right\} \; dz.
\end{align}
Computing the $z$ integral of the Green's function $G_{\mathfrak{u}} (z; \xi)$\footnote{This presumes (correctly) that the denominator in (\ref{eqn:FG_Gsfn}) does not vanish.} gives 
\begin{align}
\int_{-H(x)}^0 G_{ \mathfrak{u}}(z;\xi) \; dz & = 
\frac{- i \mathrm{e}^{-\sqrt{i} \varphi \xi} 
\left(
\mathrm{e}^{\sqrt{i} \varphi \xi} - 
\mathrm{e}^{\sqrt{i} \varphi H(x)} 
\right) \left(
\mathrm{e}^{\sqrt{i} \varphi  \left( \xi + H(x) \right)} - 1
\right)
} { \nu_0 \epsilon^2 \varphi^2  \left(1 + \mathrm{e}^{2 \sqrt{i} \varphi H(x )} \right)} , \nonumber \\
\implies \int_{-H(x)}^0 G_{ \mathfrak{u}}(z;0) \; dz & = 
\frac{ - i \left( 
\mathrm{e}^{\sqrt{i} \varphi  H(x) }   - 1   
\right)^2
} { \nu_0 \epsilon^2 \varphi^2  \left(1 + \mathrm{e}^{2 \sqrt{i} \varphi H(x )} \right)} .
\end{align}
Therefore,
\begin{align}
\mathfrak{U} (x) & = 
\frac{-i }{{ \nu_0 \epsilon^2 \varphi^2  \left(1 + \mathrm{e}^{2 \sqrt{i} \varphi H(x )} \right)} } \times \nonumber \\
&\left\{
 \int_{-H(x)}^{0} 
\mathrm{e}^{-\sqrt{i} \varphi \xi} 
\left(
\mathrm{e}^{\sqrt{i} \varphi \xi} - 
\mathrm{e}^{\sqrt{i} \varphi H(x)} 
\right) \left(
\mathrm{e}^{\sqrt{i} \varphi  \left( \xi + H(x) \right)} - 1
\right)
 \left[ \partial_x p( x, \xi) + i \partial_y p( x, \xi) \right] \; d \xi  \right. \nonumber  \\
& + \left.
\left( 
\mathrm{e}^{\sqrt{i} \varphi  H(x) }   - 1   
\right)^2
\left( \tau^x + i \tau^y \right)   \right\}
\label{eqn:UV_expression}
\end{align}
(see \texttt{PlanetaryGeostrophicTheory\_v0.2.ipynb}).

Eqns.~(\ref{eqn:U_mom}) and (\ref{eqn:V_mom}) also imply that:
\begin{align}
i f \mathfrak{U} & = - H \left(\frac{\partial p_b}{\partial x} + i \frac{\partial p_b}{\partial y} \right) - \left(\frac{\partial \chi }{\partial y} + i \frac{\partial \chi}{\partial x}  \right)  + \tau^x + i \tau^y  - \tau^x_b - i \tau^y_b .
\label{eqn:complex_vertical_int_momentum}
\end{align}
Therefore, for an unstratified fluid ($\chi = 0, \nabla p_b = \nabla p_s$), specifying the surface pressure gradient $\nabla p_s$ closes the problem. 
We compute the surface stress from (\ref{eqn:complex_vertical_int_momentum}) by substituting (\ref{eqn:bottom_stress_expression}).
Then $u(z), v(z), U, V$ and the force balance can all be computed. 
See Figure \ref{fig:example_soln} for an example calculation.
Alternatively, the surface stress can be specified, and the pressure gradient can be computed (see \texttt{Example\_barotropic\_GeostrophicFlow\_v0.3.ipynb}).
In practice, it seems that the bottom stress always vanishes. Why??  Is this correct??

\begin{figure}
\centering
\includegraphics[width=9cm]{Example_barotropic_GeostrophicFlow_v0.2.ipynb_VerticalProfiles.pdf}
\includegraphics[width=12cm]{Example_barotropic_GeostrophicFlow_v0.2.ipynb_Hodograph.pdf}
\includegraphics[width=12cm]{Example_barotropic_GeostrophicFlow_v0.2.ipynb_TermBalance.pdf}
\caption{\label{fig:example_soln} Example solution to the barotropic frictional geostrophic flow problem from \texttt{Example\_barotropic\_GeostrophicFlow\_v0.2.ipynb}.}
\end{figure}


For a stratified fluid, $b$ and hence $\chi$ (from (\ref{eqn:baroclinic_PE})) are specified. 
The total pressure field $p(x,y,z)$ consists of a known buoyancy-driven piece (baroclinic) and a surface pressure driven piece from (\ref{eqn:hydrostatic_p_pgf}).
Similarly, the bottom pressure field $p_b(x,y)$ consists of a known buoyancy-driven piece and a surface pressure driven piece from (\ref{eqn:pb}).
Thus, (\ref{eqn:complex_vertical_int_momentum}) leads to:
\begin{align}
i f \mathfrak{U}(x) & = -
\int_{-H(x)}^{0} \left( \partial_x + i \partial y \right) p(x,z) \; dz 
+ \tau_s(x)  - \tau_b (x) ,
 \end{align}
 leaving the vertically-integrated pressure gradient term explicit.
 Substituting for $\mathfrak{U}$ from (\ref{eqn:UV_expression}) and rewriting $\phi^2 = f / (\nu_0 \epsilon^2)$ gives:
 \begin{align}
& \frac{1 }{{ 1 + \mathrm{e}^{2 \sqrt{i} \varphi H(x )} } } 
\left\{
 \int_{-H(x)}^{0} 
\mathrm{e}^{-\sqrt{i} \varphi \xi} 
\left(
\mathrm{e}^{\sqrt{i} \varphi \xi} - 
\mathrm{e}^{\sqrt{i} \varphi H(x)} 
\right) \left(
\mathrm{e}^{\sqrt{i} \varphi  \left( \xi + H(x) \right)} - 1
\right)
 \left( \partial_x  + i \partial_y  \right) p( x, \xi)  \; d \xi  \right. \nonumber \\
& + \left.
\left( 
\mathrm{e}^{\sqrt{i} \varphi  H(x) }   - 1   
\right)^2
\tau_s (x)  \right\} 
= 
\tau_s (x)  - \tau_b (x) - 
\int_{-H(x)}^{0} \left( \partial_x + i \partial y \right) p(x,z) \; dz .
\end{align}
And substituting for the bottom stress $\tau_b$ from (\ref{eqn:bottom_stress_expression})
 \begin{align}
& \frac{1 }{{1 + \mathrm{e}^{2 \sqrt{i} \varphi H(x )} } } 
\left\{
 \int_{-H(x)}^{0} 
\mathrm{e}^{-\sqrt{i} \varphi \xi} 
\left(
\mathrm{e}^{\sqrt{i} \varphi \xi} - 
\mathrm{e}^{\sqrt{i} \varphi H(x)} 
\right) \left(
\mathrm{e}^{\sqrt{i} \varphi  \left( \xi + H(x) \right)} - 1
\right)
 \left( \partial_x  + i \partial_y  \right) p( x, \xi) \; d \xi  \right. \nonumber \\
& +\left.
\left( 
\mathrm{e}^{\sqrt{i} \varphi  H(x) }   - 1   
\right)^2
\tau_s (x)   \right\} 
 =
 \tau_s (x)
 - \int_{-H(x)}^{0} \left( \partial_x + i \partial y \right) p(x,z) \; dz  
 \nonumber \\
&- 
\int_{-H(x)}^{0} 
\frac{ \mathrm{e}^{-\sqrt{i} \varphi (\xi - H(x))  } }{ 1 + \mathrm{e}^{2 \sqrt{i} \varphi H \left(x \right)} }
\left( 1 + \mathrm{e}^{2 \sqrt{i} \varphi \xi}\right) 
 \left( \partial_x  + i \partial_y \right) p( x, \xi)  \; d \xi 
+\frac{ 2 \mathrm{e}^{\sqrt{i} \varphi H(x)  } }{ 1 + \mathrm{e}^{2 \sqrt{i} \varphi H \left(x \right)} }
\tau_s (x)
. \nonumber
\end{align}
Thus, substituting for the hydrostatic pressure from (\ref{eqn:hydrostatic_p}):
 \begin{align}
& \frac{1}{{ 1 + \mathrm{e}^{2 \sqrt{i} \varphi H(x )} } } \times \nonumber \\
& \left\{
 \int_{-H(x)}^{0} 
\mathrm{e}^{-\sqrt{i} \varphi \xi} 
\left(
\mathrm{e}^{\sqrt{i} \varphi \xi} - 
\mathrm{e}^{\sqrt{i} \varphi H(x)} 
\right) \left(
\mathrm{e}^{\sqrt{i} \varphi  \left( \xi + H(x) \right)} - 1
\right)
 \left( \partial_x  + i \partial_y  \right)
\left[
 p_s (x) + \int_{\xi}^{0} b(x,z') \; dz'  \right]
 \; d \xi   \right. \nonumber \\
& + \left.
\left( 
\mathrm{e}^{\sqrt{i} \varphi  H(x) }   - 1   
\right)^2
\tau_s (x)  \right\} 
 = \tau_s (x) 
 - 
 \int_{-H(x)}^{0} \left( \partial_x + i \partial y \right) \left[ 
 p_s(x,y) + \int_{z}^{0} b(x,y,z') \; dz'
 \right] \; dz  
  \nonumber \\
&- 
\int_{-H(x)}^{0} 
\frac{ \mathrm{e}^{-\sqrt{i} \varphi (\xi - H(x))  } }{ 1 + \mathrm{e}^{2 \sqrt{i} \varphi H \left(x \right)} }
\left( 1 + \mathrm{e}^{2 \sqrt{i} \varphi \xi}\right) 
  \left( \partial_x  + i \partial_y  \right)
\left[
 p_s (x) + \int_{\xi}^{0} b(x,z') \; dz'  \right]
 \; d \xi \nonumber \\
& +
\frac{ 2 \mathrm{e}^{\sqrt{i} \varphi H(x)  } }{ 1 + \mathrm{e}^{2 \sqrt{i} \varphi H \left(x \right)} }
\tau_s (x)
. \nonumber
\end{align}
Rearranging to factor $\tau_s (x)$ terms on the right gives:
 \begin{align}
& \frac{1}{{1 + \mathrm{e}^{2 \sqrt{i} \varphi H(x )} } } \times \nonumber \\
& \left\{
 \int_{-H(x)}^{0} 
\mathrm{e}^{-\sqrt{i} \varphi \xi} 
\left(
\mathrm{e}^{\sqrt{i} \varphi \xi} - 
\mathrm{e}^{\sqrt{i} \varphi H(x)} 
\right) \left(
\mathrm{e}^{\sqrt{i} \varphi  \left( \xi + H(x) \right)} - 1
\right)
 \left( \partial_x  + i \partial_y  \right)
\left[
 p_s (x) + \int_{\xi}^{0} b(x,z') \; dz'  \right]
 \; d \xi   \right\} \nonumber \\
& = \frac{4 \tau_s (x)}{\mathrm{e}^{\sqrt{i} \varphi H(x)} + \mathrm{e}^{-\sqrt{i} \varphi H(x)}}  
 - 
\int_{-H(x)}^{0} \left( \partial_x + i \partial y \right) \left[ 
 p_s(x,y) + \int_{z}^{0} b(x,y,z') \; dz'
 \right] \; dz  
  \nonumber \\
&- 
\int_{-H(x)}^{0} 
\frac{ \mathrm{e}^{-\sqrt{i} \varphi (\xi - H(x))  } }{ 1 + \mathrm{e}^{2 \sqrt{i} \varphi H \left(x \right)} }
\left( 1 + \mathrm{e}^{2 \sqrt{i} \varphi \xi}\right) 
  \left( \partial_x  + i \partial_y  \right)
\left[
 p_s (x) + \int_{\xi}^{0} b(x,z') \; dz'  \right]
 \; d \xi 
. \nonumber
\end{align}
Hence, manipulating,
 \begin{align}
& 
 \int_{-H(x)}^{0} 
\mathrm{e}^{-\sqrt{i} \varphi \xi} 
\left(
\mathrm{e}^{\sqrt{i} \varphi \xi} - 
\mathrm{e}^{\sqrt{i} \varphi H(x)} 
\right) \left(
\mathrm{e}^{\sqrt{i} \varphi  \left( \xi + H(x) \right)} - 1
\right)
 \left( \partial_x  + i \partial_y  \right)
\left[
 p_s (x) + \int_{\xi}^{0} b(x,z') \; dz'  \right]
 \; d \xi    \nonumber \\
& = 
\left(1 + \mathrm{e}^{2 \sqrt{i} \varphi H(x )} \right)
\left[ 
\frac{4 \tau_s (x)}{\mathrm{e}^{\sqrt{i} \varphi H(x)} + \mathrm{e}^{-\sqrt{i} \varphi H(x)}}  
 - 
\int_{-H(x)}^{0} \left( \partial_x + i \partial y \right) \left[ 
 p_s(x,y) + \int_{z}^{0} b(x,y,z') \; dz'
 \right] \; dz  
 \right]
\nonumber \\
&- 
\int_{-H(x)}^{0} 
\mathrm{e}^{-\sqrt{i} \varphi (\xi - H(x))  } 
\left( 1 + \mathrm{e}^{2 \sqrt{i} \varphi \xi}\right) 
  \left( \partial_x  + i \partial_y  \right)
\left[
p_s(x)  + \int_{\xi}^{0} b(x,z') \; dz'  \right]
 \; d \xi 
. \nonumber
\end{align}
Hence, breaking apart the integrals to pull the surface pressure terms out:
 \begin{align}
& 
 \left( \partial_x  + i \partial_y  \right) p_s (x) \int_{-H(x)}^{0} 
\mathrm{e}^{-\sqrt{i} \varphi \xi} 
\left(
\mathrm{e}^{\sqrt{i} \varphi \xi} - 
\mathrm{e}^{\sqrt{i} \varphi H(x)} 
\right) \left(
\mathrm{e}^{\sqrt{i} \varphi  \left( \xi + H(x) \right)} - 1
\right) \; d \xi    \nonumber \\
& +  \int_{-H(x)}^{0} 
\mathrm{e}^{-\sqrt{i} \varphi \xi} 
\left(
\mathrm{e}^{\sqrt{i} \varphi \xi} - 
\mathrm{e}^{\sqrt{i} \varphi H(x)} 
\right) \left(
\mathrm{e}^{\sqrt{i} \varphi  \left( \xi + H(x) \right)} - 1
\right)
 \left( \partial_x  + i \partial_y  \right)
 \int_{\xi}^{0} b(x,z') \; dz' 
 \; d \xi   \nonumber \\
& = 
\left(1 + \mathrm{e}^{2 \sqrt{i} \varphi H(x )} \right)
\left[ 
\frac{4 \tau_s (x)}{\mathrm{e}^{\sqrt{i} \varphi H(x)} + \mathrm{e}^{-\sqrt{i} \varphi H(x)}}  
 - 
 H(x) \left( \partial_x + i \partial y \right)  p_s(x,y) - \int_{-H(x)}^0 \left( \partial_x + i \partial y \right) \int_{z}^{0} b(x,y,z') \; dz' \; dz
\right]
\nonumber \\
&- 
  \left( \partial_x  + i \partial_y  \right)
 p_s (x) 
 \int_{-H(x)}^{0} 
\mathrm{e}^{-\sqrt{i} \varphi (\xi - H(x))  } 
\left( 1 + \mathrm{e}^{2 \sqrt{i} \varphi \xi}\right) 
 \; d \xi 
\nonumber \\
&- 
\int_{-H(x)}^{0} 
\mathrm{e}^{-\sqrt{i} \varphi (\xi - H(x))  } 
\left( 1 + \mathrm{e}^{2 \sqrt{i} \varphi \xi}\right) 
  \left( \partial_x  + i \partial_y  \right)
 \int_{\xi}^{0} b(x,z') \; dz'  
 \; d \xi 
. \nonumber
\end{align}
Moving the surface pressure terms to the left gives:
 \begin{align}
& 
 \left( \partial_x  + i \partial_y  \right) p_s (x) \left\{  \int_{-H(x)}^{0} 
\mathrm{e}^{-\sqrt{i} \varphi \xi} 
\left(
\mathrm{e}^{\sqrt{i} \varphi \xi} - 
\mathrm{e}^{\sqrt{i} \varphi H(x)} 
\right) \left(
\mathrm{e}^{\sqrt{i} \varphi  \left( \xi + H(x) \right)} - 1
\right) + \mathrm{e}^{-\sqrt{i} \varphi (\xi - H(x))  } 
\left( 1 + \mathrm{e}^{2 \sqrt{i} \varphi \xi}\right) 
\; d \xi    + \right. \nonumber \\
&\left. 
 \left(1 + \mathrm{e}^{2 \sqrt{i} \varphi H(x )} \right)
H(x) \right\} \nonumber \\
& = 
\left(1 + \mathrm{e}^{2 \sqrt{i} \varphi H(x )} \right)
\left[ 
\frac{4 \tau_s (x)}{\mathrm{e}^{\sqrt{i} \varphi H(x)} + \mathrm{e}^{-\sqrt{i} \varphi H(x)}}  
 - 
 \int_{-H(x)}^{0} \left( \partial_x + i \partial y \right) 
 \int_{z}^{0} b(x,y,z') \; dz' \; dz  
  \right]
\nonumber \\
&- 
\int_{-H(x)}^{0} 
\mathrm{e}^{-\sqrt{i} \varphi (\xi - H(x))  } 
\left( 1 + \mathrm{e}^{2 \sqrt{i} \varphi \xi}\right) 
  \left( \partial_x  + i \partial_y  \right)
 \int_{\xi}^{0} b(x,z') \; dz'  
 \; d \xi 
 \nonumber \\
 & -  \int_{-H(x)}^{0} 
\mathrm{e}^{-\sqrt{i} \varphi \xi} 
\left(
\mathrm{e}^{\sqrt{i} \varphi \xi} - 
\mathrm{e}^{\sqrt{i} \varphi H(x)} 
\right) \left(
\mathrm{e}^{\sqrt{i} \varphi  \left( \xi + H(x) \right)} - 1
\right)
 \left( \partial_x  + i \partial_y  \right)
 \int_{\xi}^{0} b(x,z') \; dz' 
 \; d \xi   \nonumber 
\end{align}
Manipulating the terms gives:
 \begin{align}
& 
 \left( \partial_x  + i \partial_y  \right) p_s (x) \left\{  \int_{-H(x)}^{0} 
\mathrm{e}^{-\sqrt{i} \varphi \xi} 
\left(
\mathrm{e}^{\sqrt{i} \varphi \xi} - 
\mathrm{e}^{\sqrt{i} \varphi H(x)} 
\right) \left(
\mathrm{e}^{\sqrt{i} \varphi  \left( \xi + H(x) \right)} - 1
\right) + \mathrm{e}^{-\sqrt{i} \varphi (\xi - H(x))  } 
\left( 1 + \mathrm{e}^{2 \sqrt{i} \varphi \xi}\right) 
\; d \xi    + \right. \nonumber \\
&\left. 
 \left(1 + \mathrm{e}^{2 \sqrt{i} \varphi H(x )} \right)
H(x) \right\} \nonumber \\
& = 
 \frac{4 \left(1 + \mathrm{e}^{2 \sqrt{i} \varphi H(x )} \right)  \tau_s (x)}{\mathrm{e}^{\sqrt{i} \varphi H(x)} + \mathrm{e}^{-\sqrt{i} \varphi H(x)}}   -  \nonumber \\
& \int_{-H(x)}^{0} 
\left[ 1 + \mathrm{e}^{2 \sqrt{i} \varphi H(x )} + \mathrm{e}^{-\sqrt{i} \varphi \xi} 
\left(
\mathrm{e}^{\sqrt{i} \varphi \xi} - 
\mathrm{e}^{\sqrt{i} \varphi H(x)} 
\right) \left(
\mathrm{e}^{\sqrt{i} \varphi  \left( \xi + H(x) \right)} - 1
\right) +
\mathrm{e}^{-\sqrt{i} \varphi (\xi - H(x))  } 
\left( 1 + \mathrm{e}^{2 \sqrt{i} \varphi \xi}\right) 
\right]
 \left( \partial_x  + i \partial_y  \right)
 \int_{\xi}^{0} b(x,z') \; dz' 
 \; d \xi   \nonumber 
\end{align}
Now define $A(x,\xi)$ as
\begin{align}
A(x,\xi) & \equiv \mathrm{e}^{-\sqrt{i} \varphi \xi} 
\left(
\mathrm{e}^{\sqrt{i} \varphi \xi} - 
\mathrm{e}^{\sqrt{i} \varphi H} 
\right) \left(
\mathrm{e}^{\sqrt{i} \varphi  \left( \xi + H \right)} - 1
\right) + \mathrm{e}^{-\sqrt{i} \varphi (\xi - H)  } 
\left( 1 + \mathrm{e}^{2 \sqrt{i} \varphi \xi}\right)   \nonumber \\
& =
\mathrm{e}^{-\sqrt{i} \varphi \xi} 
\left(
\mathrm{e}^{\sqrt{i} \varphi \left(2 \xi + H \right)} - 
\mathrm{e}^{\sqrt{i} \varphi \left(\xi + 2 H\right)}
- 
\mathrm{e}^{\sqrt{i} \varphi \xi} +
\mathrm{e}^{\sqrt{i} \varphi H} 
\right) + \mathrm{e}^{-\sqrt{i} \varphi (\xi - H)  }  + 
 \mathrm{e}^{\sqrt{i} \varphi \left( \xi + H \right)} \nonumber \\
& =
2 \mathrm{e}^{\sqrt{i} \varphi \left(\xi + H \right)} - 
\mathrm{e}^{2 \sqrt{i} \varphi  H}
- 
1 +
2 \mathrm{e}^{\sqrt{i} \varphi \left(H(x) - \xi\right)} 
\nonumber ,\\
\implies \int_{-H}^{0} A(x,\xi) \; d \xi & = 
\frac{2}{\sqrt{i} \varphi} \left( \mathrm{e}^{2 \sqrt{i} \varphi H} - 1 \right) -
 H \left( \mathrm{e}^{2 \sqrt{i} \varphi H} + 1 \right)
  \end{align}
  and
  \begin{align}
  1 + \mathrm{e}^{2 \sqrt{i} \varphi H} + A(x,\xi) = 
  2 \left( \mathrm{e}^{\sqrt{i} \varphi \left(\xi + H \right)} + \mathrm{e}^{\sqrt{i} \varphi \left(H - \xi\right)} \right) \nonumber .
  \end{align}
  Thus,
 \begin{align}
& 
 \frac{1}{\sqrt{i} \varphi} \left( \mathrm{e}^{2 \sqrt{i} \varphi H(x)} - 1 \right) 
 \left( \partial_x  + i \partial_y  \right) p_s (x) 
\nonumber \\
& = 
 \frac{2 \left(1 + \mathrm{e}^{2 \sqrt{i} \varphi H(x )} \right)  \tau_s (x)}{\mathrm{e}^{\sqrt{i} \varphi H(x)} + \mathrm{e}^{-\sqrt{i} \varphi H(x)}}
 -
\int_{-H(x)}^{0} 
\left( \mathrm{e}^{\sqrt{i} \varphi \left(H(x) - z\right)}  + \mathrm{e}^{\sqrt{i} \varphi \left(z + H(x) \right)} \right) 
 \int_{z}^{0} \left( \partial_x + i \partial y \right)  b(x,z') \; dz' \; dz     \nonumber .
\end{align}  
Hence,
 \begin{align}
& 
 \left( \partial_x  + i \partial_y  \right) p_s (x) 
\nonumber \\
& = 
 \frac{\sqrt{i} \varphi\, \tau_s (x)} { \sinh \sqrt{i} \varphi H(x)} 
-
  \frac{\sqrt{i} \varphi} { \mathrm{e}^{2 \sqrt{i} \varphi H(x)} - 1 }
\int_{-H(x)}^{0} 
\left( \mathrm{e}^{\sqrt{i} \varphi \left(H(x) - z\right)} + \mathrm{e}^{\sqrt{i} \varphi \left(z + H(x) \right)} \right) 
 \int_{z}^{0} \left( \partial_x + i \partial y \right)  b(x,z') \; dz' \; dz     \nonumber .
\end{align}  
To check this expression, see \texttt{PlanetaryGeostrophicTheory\_v0.2.ipynb}.

This final expression is of the form (now with both $(x,y)$ spatial arguments)
\begin{align}
 \left( \partial_x  + i \partial_y  \right) p_s (x,y) = F(x,y) ,
 \label{eqn:ps_diffeq}
\end{align}
where
 \begin{align}
F(x,y) = 
 \frac{\sqrt{i} \varphi\, \tau_s (x,y)} { \sinh \sqrt{i} \varphi H(x,y)} 
 -
  \frac{\sqrt{i} \varphi} { \mathrm{e}^{2 \sqrt{i} \varphi H(x,y)} - 1 }
\int_{-H(x,y)}^{0} 
\left( \mathrm{e}^{\sqrt{i} \varphi \left(H(x,y) - z\right)} + \mathrm{e}^{\sqrt{i} \varphi \left(z + H(x,y) \right)} \right) 
 \int_{z}^{0} \left( \partial_x + i \partial y \right)  b(x,y,z') \; dz' \; dz . \label{eqn:F_def}
\end{align}  
The formal solution of (\ref{eqn:ps_diffeq}) is
\begin{align}
p_s(x,y) & = \int_0^{T(x,y)} F(x -t, y - i t) \; dt ,
\end{align}
or, equivalently (because $p_s$ is real)
\begin{align}
\nabla^2 p_s(x,y) & = \Re \left[ \left( \partial_x - i \partial_y \right) F(x,y) \right] .
\end{align}
The boundary condition is that $p_s$ is a constant on the boundary and that the global integral of $p_s$ vanishes.

Notice the solvability condition derived by integrating along a closed path, i.e., 
\begin{align}
\oint_{\Gamma(x(s),y(s))}  \left( \partial_x  + i \partial_y  \right) p_s (x,y) \; ds & = \oint_{\Gamma(x(s),y(s))} F(x,y) \; ds
\end{align}
for circuit $\Gamma(x(s),y(s))$, where $s$ measures distance along the circuit.
The left hand side of this expression vanishes by the Cauchy integral theorem for holomorphic (differentiable) function $\left( \partial_x  + i \partial_y  \right) p_s (x,y)$.
Moreover, the circuit $\Gamma$ is arbitrary, which implies that 
\begin{align}
 \oint_{\Gamma(x(s),y(s))} F(x,y) \; ds & = 0 \text{~~} \forall \text{~~} \Gamma ,
 \label{eqn:RHS_line_integral}
\end{align}
which means that the circulation of $F$ along $\Gamma$ vanishes.
Therefore, the curl of $F$ must vanish everywhere:\footnote{Note Morera's theorem says that $F(x,y)$ must be holomorphic because of (\ref{eqn:RHS_line_integral}).
Therefore, $F(x,y)$ must satisfy the Cauchy-Riemann equations:
\begin{align}
\partial_x \Re \{ F \} - \partial_y \Im \{ F \} & = 0 , \nonumber \\
\partial_y \Re \{ F \} + \partial_x \Im \{ F \} & = 0 . 
\label{eqn:Cauchy_Riemann}
\end{align}
Then, writing $F$ in terms of real and imaginary parts $F = \Re \{F\} + \Im \{F\} \equiv F_1 + i F_2$, gives:
\begin{align}
\left( \partial_y  - i \partial_x\right) F &= \left( \partial_y F_1 + \partial_x F_2 \right) + i \left( \partial_y F_2 - \partial_x F_1 \right)  = 0 \nonumber 
\end{align}
from (\ref{eqn:Cauchy_Riemann}).
Notice how 
\begin{align}
& \partial_y F_1 + \partial_x F_2 \Leftrightarrow \nabla \cdot \mathbf{F} , \nonumber \\
& \partial_y F_2 - \partial_x F_1   \Leftrightarrow \nabla \times \mathbf{F} , \nonumber 
\end{align}
for vector $\mathbf{F} = (F_2; F_1)$.
 }
\begin{align}
\left( \partial_y - i \partial_x \right)F(x,y) & = 0 \text{~~} \forall \text{~~} (x,y) .
\end{align}
In words, the curl of the surface stress term in (\ref{eqn:F_def}) must equal the curl of the buoyancy gradient term in (\ref{eqn:F_def}) for the solution to satisfy the original (steady, linear) momentum equations (\ref{eqn:u_mom})--(\ref{eqn:v_mom}).


\section{Example Solutions: Work in Progress}


Next steps are to solve for:
\begin{enumerate}
%\item $b=0, H=1$ in an infinite plane.
\item $b=0, H(x,y) = 1 - x^2 - y^2$ in a circular parabolic basin. See \texttt{ExampleSolution\_v0.1.ipynb}.
\item $b(x,y)$ given by solving the vertical diffusive buoyancy equation (see section \ref{sec:b_solution}) and $H(x,y) = 1 - x^2 - y^2$ in a circular parabolic basin.
\end{enumerate}

\section{Solution to Linear Buoyancy Equation in Steady State}
\label{sec:b_solution}

At least for steady state and negligible buoyancy advection, the buoyancy equation (\ref{eqn:b}) reads:
\begin{align}
\epsilon^2 \frac{\partial }{\partial z} \left( \kappa \frac{\partial b}{\partial z} \right) + \gamma \left(  B(z) - b \right) & = 0, \label{eqn:steady_vertical_b} \\
\text{Surface~}z = 0:
b &= 0
 \\
\text{Bottom~}z = -H:
\frac{\partial b}{\partial z} &= \beta
\end{align}
for buoyancy profile $B(z)$ (in principle, $B$ could depend on $(x,y)$ too) and relaxation coefficient $\gamma$, and vertical buoyancy gradient at the sea bed $\beta$ (which may be zero). 
The internal buoyancy source allows for steady-state solutions that have non-trivial solutions (the $t\rightarrow \infty$ solution to (\ref{eqn:b})--(\ref{eqn:bot_b_bc}) is vanishing $b$ everywhere; the \citealt{peterson&callies26} solutions are transient, relaxing to zero).
This problem (\ref{eqn:steady_vertical_b}) can be solved exactly using a Green's function, at least for constant diffusivity, and I suspect also for diffusivity given by (\ref{eqn:kappa_def}). 

In detail: Solve the Green's function problem first:
\begin{align}
\frac{\epsilon^2}{\gamma} \frac{\partial }{\partial z} \left( \kappa \frac{\partial G_b}{\partial z} \right) -  G_b & =  -  \delta(z - \xi) , \label{eqn:Gsfn}\\
\text{Surface~}z = 0:
G_b &= 0 ,
\label{eqn:Gsfn_sfcbc} \\
\text{Bottom~}z = -H:
\frac{\partial G_b}{\partial z} &= 0 .
\label{eqn:Gsfn_botbc}
\end{align}
This is a second order linear ordinary differential equation with two boundary conditions, which can be solved exactly.
The solution for constant diffusivity $\kappa = \kappa_0$ is:
\begin{align}
G_b (z; \xi) & = \frac{\mathrm{e}^{\phi \left(z + \xi\right)}}
{2  \kappa_0 \phi \left(\mathrm{e}^{2 \phi H \left(x \right)} + 1\right)}
\begin{cases} 
\displaystyle \left(1 - \mathrm{e}^{- 2 \phi \xi} \right) \left(\mathrm{e}^{2 \phi H{\left(x \right)}} + \mathrm{e}^{- 2 \phi z  }\right) & \text{for}\: z \leq \xi , \\
\displaystyle \left(1 - \mathrm{e}^{- 2 \phi   z} \right) \left(\mathrm{e}^{2 \phi H{\left(x \right)}} + \mathrm{e}^{- 2 \phi \xi}\right) & \text{for}\: z \geq \xi ,
\end{cases}
\end{align}
with $\phi = \sqrt{\gamma/\kappa_0} \; \epsilon$.
Therefore, given $G_b(z; \xi)$,
\begin{align}
b(z; x) = - \int_{-H(x)}^{0} G_b(z; \xi) B(\xi; x) \; d \xi  .
\end{align}
See \texttt{DiffusionEquationSolution.ipynb} and \texttt{FindBuoyancyField\_and\_GeostrophicFlow\_v0.?.ipynb}.

In solving the momentum equations, one needs horizontal derivatives of the buoyancy field to compute the pressure gradient (\ref{eqn:hydrostatic_p_pgf}).
They are:
\begin{align}
- \partial_x b(z; x) &= \partial_x H(x)\;  G_b(z; -H(x)) \; B(-H(x); x)    \nonumber \\
&+  \int_{-H(x)}^{0} \partial_x G_b(z; \xi) \; B(\xi; x) \; d \xi  \nonumber \\
&+ \int_{-H(x)}^{0} G_b(z; \xi) \; \partial_x B(\xi; x) \; d \xi  ,
\end{align}
(and similarly for the $y$ derivative--the $y$ dependence in $b, H$, and $B$ is suppressed here for simplicity).
For constant $\kappa$, $G_b(z; -H(x))$ is a relatively straightforward expression:
\begin{align}
G_b(z; -H(x)) & = 
-(\kappa_0 \phi)^{-1} \cosh \left( \phi (z + H(x)) \right) \, \tanh \left( \phi H(x) \right) 
\end{align}
and $\partial_x G_b(z; \xi)$ is:
\begin{align}
\partial_x G_b(z; \xi) & = 
\frac{\left(- \mathrm{e}^{2 z \phi} - \mathrm{e}^{2 \xi \phi} + \mathrm{e}^{2 \phi \left(z + \xi\right)} + 1\right) \mathrm{e}^{\phi \left(- z - \xi + 2 H{\left(x \right)}\right)} }{\kappa_0 \left(\mathrm{e}^{4 \phi H{\left(x \right)}} + 2 \mathrm{e}^{2 \phi H{\left(x \right)}} + 1\right)} \partial_x H{\left(x \right)} 
\end{align}
(see \texttt{FindBuoyancyField\_and\_GeostrophicFlow\_v0.4.ipynb}.).


\bibliographystyle{agu04}
\bibliography{toms_bib}

\end{document}  